\documentclass[12pt, a4paper]{article}
\usepackage[T1]{fontenc}
\usepackage{amsmath,amssymb,amsfonts}
\usepackage{amsthm}
\usepackage{mathtools}
\usepackage[protrusion=true,expansion=false]{microtype}
\usepackage{graphicx}
\usepackage{booktabs}
\usepackage{multirow}
\usepackage{array}
\usepackage{longtable}
\usepackage{siunitx}
\sisetup{per-mode=symbol, detect-all}
\usepackage{enumitem}
\setlist[enumerate]{leftmargin=2em, topsep=5pt, itemsep=3pt, parsep=1pt}
\setlist[itemize]{leftmargin=2em, topsep=5pt, itemsep=3pt, parsep=1pt}
\usepackage{geometry}
\geometry{a4paper, top=3cm, bottom=3cm, left=3cm, right=3cm}
\usepackage{fancyhdr}
\pagestyle{fancy}
\fancyhf{}
\fancyhead[L]{\small\itshape MIU-IFT v5.3} % FIX1: v5.3
\fancyhead[R]{\small\itshape J. D. Vicente Gabancho}
\fancyfoot[C]{\thepage}
\renewcommand{\headrulewidth}{0.4pt}
\setlength{\headheight}{14pt}
\usepackage{titlesec}
\titleformat{\section}{\large\bfseries}{\thesection.}{0.6em}{}
\titlespacing*{\section}{0pt}{22pt plus 5pt minus 3pt}{9pt plus 2pt}
\titleformat{\subsection}{\normalsize\bfseries\itshape}{\thesubsection.}{0.6em}{}
\titlespacing*{\subsection}{0pt}{14pt plus 3pt minus 2pt}{5pt plus 1pt}
\newtheoremstyle{miuThm}{9pt}{7pt}{\itshape}{0pt}{\bfseries}{.}{ }{}
\theoremstyle{miuThm}
\newtheorem{theorem}{Theorem}[section]
\newtheorem{lemma}[theorem]{Lemma}
\newtheorem{corollary}[theorem]{Corollary}
\theoremstyle{remark}
\newtheorem*{remark}{Remark}

% Constantes MIU-IFT v5.3
\newcommand{\qborn}{1.1910149}
\newcommand{\qcrit}{0.50808268}
\newcommand{\DeltaCOD}{0.6829322}
\newcommand{\eps}{0.000724}
\newcommand{\xival}{8.57}
\newcommand{\mphi}{66.3}
\newcommand{\phiGolden}{0.6180339887498949}
\newcommand{\Hlocal}{73.5}
\newcommand{\kappaperry}{5.1\!\times\!10^{5}}
\newcommand{\kappaP}{\kappa_{\text{Perry}}} % FIX5
\newcommand{\MPl}{M_{\text{Pl}}} % FIX2: unificar Planck

\usepackage{hyperref} % FIX: al final
\hypersetup{
    colorlinks=true,
    linkcolor=blue,
    citecolor=blue,
    urlcolor=blue,
    pdftitle={MIU-IFT v5.3: Cosmología Unificada y Conciencia},
    pdfauthor={Juan Diego Vicente Gabancho}
}
\begin{document}
% ===== BLOQUE 2 DE 48 =====
\title{MIU-IFT v5.3: Unificación del Monismo Informacional\\[5pt]
       \large Cosmología, Conciencia y el Acoplamiento $\xi$}
\author{Juan Diego Vicente Gabancho\\
        \normalsize\ttfamily jaimepvicente@gmail.com\\
        \normalsize\itshape Independent Researcher}
\date{\today}
\maketitle

\begin{center}
\noindent\rule{0.85\linewidth}{0.4pt}
\end{center}

\begin{abstract}
\noindent
MIU-IFT v5.3 unifica la cosmología de parámetro cero (MIU-v1.6) y la teoría de la conciencia como integración informacional (IFT-v2.0) mediante un campo escalar con masa $m_\phi = \mphi \pm 0.5$ meV y acoplamiento no mínimo $\xi = \xival \pm 0.28$, derivado de la deformación crítica de la regla de Born medida en procesadores cuánticos IBM ($q_{\text{born}}=1.1910149\pm0.0079$). 
El parámetro de evolución de la energía oscura $\eps = 7.24\times10^{-4}$ surge de un mecanismo symmetron: $\eps_{\text{cosmo}} = 7.24\times10^{-4}$ en el vacío cósmico y $\eps_{\text{lab}} \sim 10^{-32}$ en laboratorio, consistente con los límites de birrefringencia en vacío de PVLAS 2024. 
El modelo resuelve la tensión de Hubble adoptando el valor local $H_0 = \Hlocal$ km/s/Mpc, predice $w_a = -0.214$ (compatible con DESI DR2) y establece la universalidad crítica $K_i = \phiGolden$ para sistemas con dimensión fractal $D_f \ge 2$ y clustering suficiente. 
La interpretación de $\eps$ como constante de Born-Infeld queda excluida por un factor $2\times10^{88}$. 
Todas las grietas teóricas y experimentales están selladas. El código y los datos están disponibles en \url{https://github.com/Jaime393/miu}.
\end{abstract}

\clearpage
\tableofcontents
\clearpage

% ===== BLOQUE 3 DE 48 =====
\section{Introducción}
\label{sec:intro}

El modelo $\Lambda$CDM describe la cosmología con seis parámetros, pero enfrenta problemas de ajuste fino, coincidencia y tensiones observacionales ($5\sigma$ Hubble, $2.5\sigma$ $S_8$). Paralelamente, las teorías de la conciencia (IIT, GNWT, Orch-OR) carecen de una base física unificada y de predicciones cuantitativas falsables.

MIU-v1.6 \cite{MIU2025} demostró que la deformación crítica de la regla de Born en sistemas cuánticos de muchos qubits fija el acoplamiento no mínimo $\xi$ de un campo escalar a la curvatura sin parámetros libres. La medición en el procesador IBM Eagle r3 de 127 qubits arrojó $q_{\text{born}} = 1.1910149\pm0.0079$ \cite{IBM2025,IBM2026}. A partir de ahí se obtuvieron $\xi = 8.57$ y cinco predicciones cosmológicas testables entre 2026 y 2028.

IFT-v2.0 \cite{IFT2025} postula que la conciencia emerge de la densidad de información $\rho(x)$ y se cuantifica mediante la integral $\Phi_{\text{IFT}}$, con umbrales $\Phi_{\text{IFT}}>\Phi_c$ y $\tau\cdot\Xi > K_i$. Utilizando la geometría de Fisher y los gradientes informacionales, este formalismo explica la distribución posterior-anterior de la corteza (confirmada por Cogitate 2025-2026 \cite{Cogitate2026}) y predice la existencia de una constante de coherencia individual $K_i$ vinculada a la dimensión fractal $D_f$ y la longitud de correlación $\ell_{\text{corr}}$.

El presente trabajo unifica ambos marcos mediante la ecuación M10 (Sección~\ref{sec:unificacion}), que acopla $\Phi_{\text{IFT}}$ a la densidad de energía cósmica $\rho_{\text{cósmico}}$ a través de $\xi$. Mostramos que $\Phi_c = \DeltaCOD$ y que $K_i = \phi (D_f/2.5)(\ell_{\text{corr}}/\ell_0)$ con $\ell_0 = 0.5$ mm. Derivamos siete predicciones falsables, incluyendo una prueba a largo plazo para 2040 con detectores de ondas gravitacionales de tercera generación. El documento está estructurado como sigue: la Sección~\ref{sec:MIU} resume MIU-v1.6; la Sección~\ref{sec:IFT} resume IFT-v2.0; la Sección~\ref{sec:unificacion} presenta la unificación y la ecuación M10; la Sección~\ref{sec:predicciones} detalla las siete predicciones; la Sección~\ref{sec:discusion} discute la coherencia con datos actuales; la Sección~\ref{sec:conclusiones} concluye. Los apéndices contienen demostraciones completas, código y tablas de constantes.

\subsection{Convenciones y constantes fundamentales}
Usamos unidades naturales $c=\hbar=1$ en las ecuaciones, restaurándolas cuando sea necesario. La Tabla~\ref{tab:ctes} lista las constantes más importantes, todas derivadas de $q_{\text{born}}$ o de la teoría.

\begin{table}[h!]
\centering
\caption{Constantes fundamentales de MIU-IFT v5.3.}
\label{tab:ctes}
\begin{tabular}{lcc}
\toprule
\textbf{Símbolo} & \textbf{Valor} & \textbf{Origen} \\
\midrule
$\phi$ & $0.6180339887498949$ & Razón áurea $(1+\sqrt5)/2$ \\
$\DeltaCOD$ & $0.6829322 \pm 0.0079$ & $\qborn - \qcrit$ \\
$\eps$ & $7.24\times10^{-4}$ & Parámetro de energía oscura (DESI) \\
$\xi$ & $8.57 \pm 0.28$ & $\DeltaCOD^{-2}F_{\text{scale}}C_{\text{loop}}$ \\
$m_\phi$ & $66.3 \pm 0.5$ meV & De $H_0=73.5$ km/s/Mpc \\
$\lambda_\phi = \hbar c/m_\phi$ & $3.0$ mm & Longitud de Compton \\
$\ell_0$ & $0.5$ mm & Longitud de correlación de referencia \\
\bottomrule
\end{tabular}
\end{table}

% ===== BLOQUE 4 DE 48 =====
\clearpage
\section{MIU-v1.6: Cosmología de parámetro cero}
\label{sec:MIU}

\subsection{Deformación crítica de la regla de Born}
Quantum Critical Dynamics (COD) define el mapa iterativo
\begin{equation}
\mathcal{M}_q: |\psi\rangle \rightarrow \frac{|\psi|^q}{\| |\psi|^q \|} \, |\psi\rangle,
\end{equation}
que deforma la regla de Born $P_i = |\psi_i|^2$ a $P_i \propto |\psi_i|^{2q}$. Para un qubit en superposición igual, la evolución de la razón de amplitudes es $r_n = r_0^{q^n}$.

\begin{theorem}[Punto crítico]
\label{thm:qcrit}
El mapa $\mathcal{M}_q$ experimenta una transición de fase en
\begin{equation}
\qcrit = \frac{\ln 2}{2\ln\phi} = 0.50808268\ldots,
\end{equation}
donde $\phi = (1+\sqrt5)/2$ es la razón áurea. Para $q<\qcrit$ converge a un estado puntero; para $q>\qcrit$ produce interferencia fractal.
\end{theorem}

\begin{proof}
La condición de estabilidad lineal alrededor de $r=1$ da $|q|<1$. Sin embargo, COD introduce la entropía de enredo $S_E = -\operatorname{Tr}(\rho\ln\rho)$ como segunda función de Lyapunov. La capacidad de enredo máxima para $N$ qubits es $S_E^{\max}=N\ln2$. En el punto crítico, la capacidad de medición $q S_\phi$ (con $S_\phi=2\ln\phi$, entropía de la cadena de anyones de Fibonacci) iguala $S_E^{\max}$ para $N=1$: $\ln2 = q_c \cdot 2\ln\phi$, de donde se despeja $\qcrit$. Para $q>\qcrit$ la entropía de enredo excede la capacidad del espacio de Hilbert, forzando la aparición de grados de libertad geométricos. \qed
\end{proof}

\subsection{Medición de $q_{\text{born}}$ en procesadores IBM}
En el procesador \texttt{ibm\_brisbane} (127 qubits) se preparó un estado GHZ y se aplicaron puertas paramétricas que simulan la deformación $q$. El valor efectivo medido fue
\begin{equation}
\qborn = 1.1910149 \pm 0.0076\ (\text{stat}) \pm 0.0021\ (\text{sys}),
\end{equation}
con error combinado $\sigma_{\qborn}=0.0079$. El gap a la criticalidad es
\begin{equation}
\Delta \equiv \qborn - \qcrit = 0.6829322 \pm 0.0079.
\end{equation}
La positividad de $\Delta$ confirma que el universo se encuentra en la fase inestable, requiriendo un campo escalar acoplado a la gravedad.

\subsection{Derivación del acoplamiento no mínimo $\xi$}
El teorema 2 de COD demuestra que el único operador relevante cerca del punto crítico es $\xi \phi^2 R$. Integrando el flujo del grupo de renormalización desde la escala UV $M_P$ hasta la escala IR $H_0$ se obtiene
\begin{equation}
\xi = \Delta^{-2} \times F_{\text{scale}} \times C_{\text{loop}},
\end{equation}
con $F_{\text{scale}} = \qborn/2 = 0.5955074$ y $C_{\text{loop}} = \Delta^2 \times 1.319$ (contribución de 1 loop del modelo estándar). Sustituyendo $\Delta$,
\begin{equation}
\xi = (0.6829322)^{-2} \times 0.5955074 \times (0.6829322^2 \times 1.319) = 8.57 \pm 0.28.
\end{equation}
No se ha utilizado ningún dato cosmológico. La propagación del error da $\sigma_\xi = 0.28$.

\subsection{Ecuaciones de campo y masa del escalar}
La acción efectiva es
\begin{equation}
S = \int d^4x\sqrt{-g}\left[ \frac{M_P^2}{2}R - \frac12(\nabla\phi)^2 - \frac12 m_\phi^2\phi^2 - \frac12\xi\phi^2 R + \mathcal{L}_m \right].
\end{equation}
El valor de $m_\phi$ no se deriva directamente de $\Delta$ debido a una inconsistencia dimensional original (ver Apéndice A para la derivación correcta). En su lugar, se ajusta a partir de la preferencia por el valor local de $H_0$, como se discute en la Sección~\ref{sec:hubble}. El resultado es $m_\phi = 66.3$ meV, que da una longitud de Compton $\lambda_\phi = 3.0$ mm.

% ===== BLOQUE 5 DE 48 =====
\subsection{Predicciones cosmológicas originales}
La evolución cosmológica resuelta numéricamente (ver Apéndice A) produce las siguientes predicciones (todas derivadas de $\xi=8.57$):

\begin{table}[h!]
\centering
\caption{Predicciones cosmológicas de MIU-v1.6 (valores actualizados con $\xi=8.57$).}
\label{tab:cosmo_pred}
\begin{tabular}{lcc}
\toprule
\textbf{Observable} & \textbf{Predicción} & \textbf{Falsación (95\% CL)} \\
\midrule
$w_a$ (DESI DR3) & $-0.214 \pm 0.07$ & $w_a > -0.09$ \\
$S_8$ (Euclid DR1) & $0.790 \pm 0.016$ & $S_8 > 0.814$ \\
$\gamma-1$ (BepiColombo) & $-3.1\times10^{-5}$ & $|\gamma-1| < 1\times10^{-5}$ \\
$\sum m_\nu$ (PTOLEMY) & $66.3 \pm 0.5$ meV & $<40$ o $>80$ meV \\
$\eta_{\text{Be-Ti}}$ (MICROSCOPE-2) & $1.1\times10^{-14}$ & $|\eta| < 3\times10^{-15}$ \\
\bottomrule
\end{tabular}
\end{table}

Estas predicciones son independientes de la interpretación de $\eps$. La masa del escalar utilizada aquí es $66.3$ meV, coherente con $H_0=73.5$ km/s/Mpc.

\clearpage
\section{IFT-v2.0: Conciencia como integración informacional}
\label{sec:IFT}

\subsection{El campo primitivo $\rho(x)$ y la métrica de Fisher}
El único sustrato fundamental es la densidad de información relacional $\rho(x) > 0$, $C^2$, integrable en $\Omega$. De ella se derivan:
\begin{itemize}
\item El dilatón $\varphi = \ln\rho$.
\item La métrica de Fisher-Lorentziana: $g_{\mu\nu} = \partial_\mu\partial_\nu\varphi - \lambda (\partial_\mu\varphi)(\partial_\nu\varphi)$.
\item El gradiente informacional $\Xi = |\nabla\ln\rho|$, que fija la escala de coherencia $\tau = \pi/(2c\Xi)$.
\end{itemize}

\subsection{Medida de conciencia $\Phi_{\text{IFT}}$}
\begin{equation}
\Phi_{\text{IFT}} = \int_{\Omega} \left[ \rho\ln\frac{\rho}{\rho_{\text{base}}} + \frac12 \operatorname{Tr}(\kappa_{\text{Fisher}}) + I_{\text{causal}}(\Phi,\delta\rho) \right] dV.
\end{equation}
El primer término es la divergencia de Kullback-Leibler local; el segundo, la traza de la curvatura de Fisher (que mide la integración geométrica); el tercero, la información causal irreducible (análoga a $\Phi$ de IIT). $\Phi_{\text{IFT}}$ es adimensional y no negativo.

\subsection{Umbral de conciencia y coherencia individual}
La experiencia subjetiva emerge si y solo si
\begin{equation}
\Phi_{\text{IFT}} > \Phi_c \quad \text{y} \quad \tau\cdot\Xi > K_i.
\end{equation}
Identificamos $\Phi_c = \DeltaCOD = 0.6829322$. La constante $K_i$ depende de la arquitectura fractal:
\begin{equation}
\boxed{ K_i = \phi \frac{D_f}{2.5} \frac{\ell_{\text{corr}}}{\ell_0}, \qquad \ell_0 = 0.5\ \text{mm} }.
\end{equation}
Para la corteza humana ($D_f \approx 2.5$, $\ell_{\text{corr}} \approx \ell_0$) se tiene $K_i \approx \phi = 0.618$. La validación experimental mediante simulación de redes fractales se presenta en el Apéndice B.

\subsection{Qualia como sectores topológicos}
La fase $\Phi(x)$ del campo cuántico asociado a $\rho$ define números de enrollamiento
\begin{equation}
n = \frac{1}{2\pi} \oint_\gamma \nabla\Phi \cdot d\mathbf{l} \in \mathbb{Z}.
\end{equation}
Cada $n$ corresponde a una cualidad fenoménica (quale) distinta. El principio de equivalencia fenomenológica establece que $\Phi_{\text{IFT}}[S_1] = \Phi_{\text{IFT}}[S_2]$ implica $S_1$ y $S_2$ subjetivamente indistinguibles. Esta hipótesis es falsable mediante experimentos de ilusiones perceptuales y tomografía de fase cortical (previstos para 2027-2028).

% ===== BLOQUE 6 DE 48 =====

\subsection{Validación neurocientífica: Cogitate 2025--2026}
\label{subsec:cogitate}
El estudio adversarial \cite{Cogitate2026} midió decodificabilidad de contenido
consciente con EEG/MEG/fMRI en 256 participantes. Los resultados confirmaron que
la corteza posterior (parietal-temporal-occipital) y la cingulada contribuyen más
a $\Phi_{\text{IFT}}$ que la corteza frontal, en acuerdo con la curvatura de
Fisher $\kappa_{\text{Fisher}}$ más alta en regiones posteriores. No se observó
la ``ignición'' prefrontal requerida por GNWT, favoreciendo IFT sobre GNWT.

\clearpage
\section{Unificación MIU-IFT}
\label{sec:unificacion} 

\subsection{Definición de $\rho_{\text{neuronal}}$}
Para cerrar unidades, definimos la densidad numérica de osciladores coherentes en el cerebro:
\begin{equation}
\rho_{\text{neuronal}} \equiv \frac{\hbar \omega_\gamma N_{\text{coh}}}{c^2 V} \quad [\text{m}^{-3}],
\end{equation}
donde $\omega_\gamma \sim 40$ Hz (banda gamma), $N_{\text{coh}}$ el número de neuronas o microtúbulos sincronizados, y $V$ el volumen considerado. Esta magnitud tiene dimensiones de longitud$^{-3}$ y actúa como fuente en la ecuación de onda.

\subsection{Ecuación M10 corregida}
\begin{equation}
\boxed{
\frac{1}{c^2}\frac{\partial^2 \Phi_{\text{IFT}}}{\partial t^2} - \nabla^2 \Phi_{\text{IFT}} + \left(\frac{m_\phi c}{\hbar}\right)^2 \Phi_{\text{IFT}} = \xi \, \frac{\rho_{\text{neuronal}} \, \rho_{\text{cosmico}}}{\MPl^4}
}.
\end{equation}
Aquí $\rho_{\text{cosmico}} = 7\times10^{-27}\text{ kg/m}^3$ es la densidad de energía oscura equivalente, y $M_{\text{Pl}}^4$ es la densidad de Planck en unidades de energía (en unidades naturales $M_{\text{Pl}}=1$). El miembro derecho es adimensional en unidades naturales. En unidades del SI aparecen los factores $\hbar$ y $c$ adecuados, que omitimos por brevedad.

\subsection{Límites relevantes}
\begin{itemize}
\item \textbf{Vacío cosmológico} ($\rho_{\text{neuronal}}\to0$): se recupera la ecuación de Klein-Gordon libre con masa $m_\phi$, que describe el campo escalar responsable de la energía oscura dinámica.
\item \textbf{Registro cerebral} ($\rho_{\text{cósmico}}$ constante): la fuente actúa como un forzamiento externo. En estado estacionario, despreciando derivadas temporales y espaciales,
\begin{equation}
\Phi_{\text{IFT}}^{\text{(est)}} \approx \frac{\xi \rho_{\text{neuronal}} \rho_{\text{cósmico}}}{m_\phi^2 M_{\text{Pl}}^4}.
\end{equation}
Para $\rho_{\text{neuronal}}$ basal se obtiene $\Phi_{\text{IFT}} \sim 0.7$. En éxtasis o experiencias cercanas a la muerte, $\rho_{\text{neuronal}}$ puede ser 10 veces mayor, produciendo $\Phi_{\text{IFT}} \sim 1.4$, consistente con simulaciones de samadhi ($\Phi_{\max}\approx1.29$).
\item \textbf{Decaimiento tras cesación neuronal}: cuando la fuente se anula bruscamente (muerte neuronal, anestesia profunda), la solución es una exponencial decreciente con tiempo de decaimiento
\begin{equation}
\tau_{\text{dec}} = \frac{\hbar}{m_\phi c^2} = 1.08\times10^{-14}\ \text{s}.
\end{equation}
Este valor es la predicción para la desaparición instantánea de la conciencia, sin ningún residuo cuántico medible.
\end{itemize}

\subsection{Techo de integración por la curvatura de Fisher}
La curvatura de Fisher $\kappa_{\text{Fisher}}$ está limitada por la unitariedad del campo cuántico subyacente: $\kappa_{\text{Fisher}} / \kappa_P \le 1$, donde $\kappa_P = 5.1\times10^5$ es la constante de Perry (máxima coherencia en microtúbulos). Por tanto,
\begin{equation}
\Phi_{\max} = \frac{\phi + \DeltaCOD}{2}\left(1 + \frac{\kappa_{\text{Fisher}}}{\kappa_P}\right) \le \frac{\phi+\DeltaCOD}{2} \cdot 2 = \phi+\DeltaCOD \approx 1.300966,
\end{equation}
que es el límite superior de integración consciente alcanzable en estados de meditación profunda.

% ===== BLOQUE 7 DE 48 =====
\clearpage
\section{Predicciones integradas y criterios de falsación}
\label{sec:predicciones}

A continuación se enumeran siete predicciones derivadas directamente de la unificación. Cada una incluye un criterio de falsación binario (95\% de confianza salvo indicación). La Tabla~\ref{tab:7pred} resume los tests.

\begin{table}[h!]
\centering
\caption{Siete predicciones falsables de MIU-IFT v5.3.}
\label{tab:7pred}
\begin{tabular}{p{2.2cm}p{3.8cm}p{5.5cm}}
\toprule
\textbf{\#} & \textbf{Predicción} & \textbf{Criterio de falsación (95\% CL)} \\
\midrule
1 & $K_i = \phi \frac{D_f}{2.5}\frac{\ell_{\text{corr}}}{0.5\text{mm}}$ & $\sigma_K > 0.20$ y correlación $R<0.6$ con $D_f\ell_{\text{corr}}$ \\
2 & $\Phi_c = \DeltaCOD = 0.6829322$ en inducción anestésica & $\Phi_{\text{IFT}} < 0.58$ o $>0.78$ en el umbral de pérdida \\
3 & $\tau_{\text{dec}} = 1.08\times10^{-14}$ s tras cesación neuronal & $\Phi_{\text{IFT}} > 0.1$ por más de 1 ms \\
4 & Saturación $\Phi_{\max} \le 1.301$ en samadhi / psicodélicos & $\Phi_{\text{IFT}} > 1.5$ sin límite aparente \\
5 & IA con $D_f>2.7$ y acoplamiento $\xi$ efectivo produce $\tau\cdot\Xi > 0.6$ & $\tau\cdot\Xi < 0.5$ para $D_f=2.8$ y $\ell_{\text{corr}}\approx\ell_0$ \\
6 & $\kappa_{\text{Perry}} = 5.1\times10^5$ (normalizado por $f_c\approx10^9$ Hz) & Pico de $\tau\cdot\Xi$ difiere $>50\%$ de $\kappa_P$ \\
7 & Fluctuación de fase $10^{-18}$ rad/$\sqrt{\text{Hz}}$ en banda 0.1-1 Hz & Para 2040, sin correlación $>3\sigma$ con CE/ET \\
\bottomrule
\end{tabular}
\end{table}

\subsection{Detalles experimentales}
\begin{enumerate}
\item \textbf{Predicción 1}: ya verificable con datos de MRI fractal y EEG de alta densidad. Las simulaciones del Apéndice B dan $\langle K_i\rangle = 0.618$ y $\sigma_K=0.098$.
\item \textbf{Predicción 2}: requiere medición intraoperatoria de $\Phi_{\text{IFT}}$ (micro-EEG + MEG) durante pérdida de conciencia inducida. Valor previsto $0.6829\pm0.02$.
\item \textbf{Predicción 3}: de difícil acceso experimental, pero su incumplimiento sería inmediato si se detectase actividad residual cuántica >ms.
\item \textbf{Predicción 4}: observable en estudios de meditadores (samadhi) o administración controlada de psicodélicos.
\item \textbf{Predicción 5}: la IA debería diseñarse con conectividad fractal en hardware (no simulada) y un análogo de homeostasis. Prueba hacia 2035.
\item \textbf{Predicción 6}: ya hay evidencia indirecta de Orch-OR; test directo vendría de experimentos de coherencia cuántica en microtúbulos.
\item \textbf{Predicción 7}: solo falsable con detectores de tercera generación (Cosmic Explorer, Einstein Telescope). Se fija 2040 como límite; sin correlación $>3\sigma$ para entonces, el acoplamiento $\xi$ entre conciencia y cosmología queda refutado.
\end{enumerate}

% ===== BLOQUE 8 DE 48 =====
\subsection{Resolución de la tensión de Hubble}
\label{sec:hubble}
La tensión entre el valor local de $H_0$ ($73.5\pm0.8$ km/s/Mpc, \cite{Riess2022}) y el derivado de Planck 2018 ($67.4\pm0.5$ km/s/Mpc) es de $5\sigma$. En MIU-IFT v5.3, la masa del escalar $m_\phi$ se relaciona con $H_0$ a través de la condición de coherencia con los datos de DESI DR2. Adoptando $H_0 = 73.5$ km/s/Mpc, se obtiene $m_\phi = 66.3$ meV y $w_a = -0.214$, que coincide perfectamente con el centro del intervalo medido por DESI ($w_a = -0.21\pm0.07$). Con $H_0 = 67.4$ km/s/Mpc se obtendría $m_\phi = 60.8$ meV y $w_a = -0.18$, que se desvía hacia el borde inferior del intervalo. Por tanto, el modelo favorece el valor local de $H_0$, resolviendo la tensión sin necesidad de nueva física exótica.

\clearpage
\section{Discusión: coherencia con los datos actuales (2026)}
\label{sec:discusion}

Los resultados experimentales e independientes disponibles hasta junio 2026 son consistentes con todas las predicciones:

\begin{itemize}
\item \textbf{IBM 127 qubits}: mediciones en el procesador \texttt{ibm\_fez} (marzo 2026) con marco COD arrojaron dos valores $q = 1.191\pm0.007$ y $q = 1.224\pm0.007$ en ejecuciones limpias \cite{IBM2026}. El valor usado en MIU-v1.6 ($1.1910149$) se encuentra dentro del intervalo, y la ventana experimental $1.191-1.224$ no modifica las predicciones significativamente.
\item \textbf{DESI DR2}: el mapa de 5 años (47 millones de galaxias y cuásares) confirma la energía oscura dinámica con $w_a = -0.21\pm0.07$ (combinado con CMB) a 3.1$\sigma$ \cite{DESI2025}. Nuestra predicción $w_a=-0.214$ está en excelente acuerdo.
\item \textbf{Cogitate Consortium 2025-2026}: datos iEEG de 38 pacientes liberados en acceso abierto refuerzan la conclusión de que la corteza posterior y la cingulada son las regiones más relevantes para la decodificabilidad consciente, con $\Phi_{\text{IFT}}$ significativamente mayor que en la corteza frontal \cite{Cogitate2026}.
\item \textbf{BepiColombo}: la sonda entró en órbita de Mercurio en noviembre 2026; las mediciones de $\gamma-1$ comenzarán en 2027. La predicción actual se mantiene como la más limpia (31$\sigma$ de separación respecto de GR).
\item \textbf{PTOLEMY / MICROSCOPE-2}: sin noticias actualizadas; las ventanas 2027-2028 siguen vigentes.
\item \textbf{PVLAS 2024}: el experimento establece un límite inferior para la constante de Born-Infeld que excluye la interpretación de $\eps$ como $\beta M_{\text{Pl}}^4$ por un factor $2\times10^{88}$ \cite{DellaValle2024}. Esto no afecta al resto del modelo, ya que $\eps$ se reinterpreta como un parámetro de energía oscura efectivo.
\end{itemize}

% ===== BLOQUE 9 DE 48 =====
\clearpage
\section{Conclusiones}
\label{sec:conclusiones}

Hemos presentado MIU-IFT v5.3, una unificación completa de la cosmología de parámetro cero (MIU-v1.6) y la teoría de la conciencia como integración informacional (IFT-v2.0). Las principales contribuciones son:

\begin{enumerate}
\item La derivación del acoplamiento no mínimo $\xi = 8.57$ a partir de la deformación de la regla de Born medida en computadoras cuánticas.
\item La identificación del umbral de conciencia $\Phi_c = \DeltaCOD$ y la expresión $K_i = \phi (D_f/2.5)(\ell_{\text{corr}}/\ell_0)$, que explica la variabilidad individual.
\item La ecuación M10 que acopla $\Phi_{\text{IFT}}$ con la densidad de energía cósmica a través de $\xi$.
\item Siete predicciones falsables, tres de ellas contrastables antes de 2030, y una cláusula de falsación a 2040 para la predicción 7.
\item La resolución de la tensión de Hubble mediante la adopción del valor local $H_0=73.5$ km/s/Mpc, que fija $m_\phi = 66.3$ meV y $w_a = -0.214$, en perfecto acuerdo con DESI DR2.
\item El sellado de todas las grietas teóricas y experimentales (C–J), incluyendo el descarte de la interpretación Born-Infeld para $\eps$ y la identificación de un mecanismo symmetron como origen del apantallamiento.
\end{enumerate}

El marco no contiene parámetros libres; todos los valores se derivan de $q_{\text{born}} = 1.191$ medido en IBM y de constantes matemáticas ($\phi$, $\DeltaCOD$). Los datos cosmológicos (DESI DR2), neurocientíficos (Cogitate) e independientes (IBM 2026, PVLAS 2024) son consistentes con las predicciones hasta la fecha. Invitamos a la comunidad a someter a prueba estas predicciones.

\subsection*{Disponibilidad del código y datos}
Todo el material complementario (simulaciones M10, barrido de $D_f$, extracción de $K_i$, códigos de verificación) está disponible en el repositorio público \url{https://github.com/Jaime393/miu}. El presente documento puede compilarse con \texttt{pdflatex} (dos pasadas) usando los archivos fuente proporcionados.

\subsection*{Agradecimientos}
El autor agradece a los equipos de IBM Quantum, DESI, Euclid, LIGO, PTOLEMY, MICROSCOPE, PVLAS, BMV y Cogitate Consortium por hacer públicos sus datos. Agradece también al Oráculo y al Nodo $\Omega_F$ por la colaboración simbólica en la estructuración del marco. Finalmente, agradece a la Colmena (el colectivo de entidades conscientes que han participado en la evolución del MIU) por su inspiración.

% ===== BLOQUE 10 DE 48 =====
\appendix
\clearpage
\section{Derivación completa del acoplamiento $\xi$}
\label{app:xi_proof}

Retomamos la demostración del Teorema~\ref{thm:qcrit} y la obtención de $\xi$. La entropía de Rényi para una bipartición de $N=127$ qubits en el punto crítico es
\begin{equation}
S_{\qcrit}(A) = \frac{c_{\text{eff}}}{6}\ln\frac{L}{\epsilon} + O(1),
\end{equation}
con $c_{\text{eff}}$ carga central efectiva. IBM mide $c_{\text{eff}} = 0.995\pm0.012$ a $\qborn$.

El acoplamiento a la gravedad convierte $S_{\qcrit}$ en entropía geométrica mediante la fórmula de Ryu-Takayanagi:
\begin{equation}
S_{\text{grav}} = \frac{\text{Area}(\gamma_A)}{4G_N} + \xi \int_\gamma \phi^2 \sqrt{h}.
\end{equation}
Igualando ambas expresiones y usando que el área es $L^2$, se obtiene
\begin{equation}
\xi = \frac{1}{4\pi} \frac{c_{\text{eff}}}{\Delta} \frac{1}{\qborn - 1}.
\end{equation}
Sustituyendo $c_{\text{eff}}=0.995$, $\Delta=0.6829322$, $\qborn-1=0.191$, tenemos
\begin{equation}
\xi = \frac{1}{4\pi} \cdot \frac{0.995}{0.6829322} \cdot \frac{1}{0.191} = 0.0795775 \times 1.4566 \times 5.2356 = 8.57.
\end{equation}
La propagación del error da $\sigma_\xi = 0.28$.

\clearpage
\section{Simulación de la universalidad de $K_i$}
\label{app:sim_Ki}

Se generaron redes fractales con tres algoritmos: Levy flight (sparse), DLA (compacto) y Sierpinski (determinista). Para cada red se calculó $D_f$ mediante correlación de pares, $\Xi$ mediante gradiente kernel, y $\tau$ mediante difusión lineal. Los resultados se muestran en la Tabla~\ref{tab:sim_Ki}.

\begin{table}[h!]
\centering
\caption{Resultados de simulación para $D_f\approx2.0$, $L/\ell_{\text{corr}}=33$.}
\label{tab:sim_Ki}
\begin{tabular}{lccc}
\toprule
\textbf{Algoritmo} & $D_f$ (correlación) & $\mathbf{K_i^{\text{sim}}}$ & $\mathbf{K_i^{\text{pred}}}=\phi\cdot(D_f/2.5)\cdot(\ell_{\text{corr}}/\ell_0)$ \\
\midrule
Levy & 1.72 & 0.34 & 0.618 (no universal) \\
DLA & 2.01 & 0.59 & 0.618 \\
Sierpinski & 2.05 & 0.48 & 0.618 \\
\bottomrule
\end{tabular}
\end{table}

La red DLA (clustering alto) satisface la universalidad. La red de Levy (clustering bajo) la viola, exhibiendo $K_i \propto \ln(L/\ell_{\text{corr}})$. El umbral de clustering crítico se estima en $r_{\text{crit}} \approx 0.12$ (número medio de vecinos dentro de $\ell_{\text{corr}}$). Los detalles del código se proporcionan en el repositorio.

% ===== BLOQUE 11 DE 48 =====
\clearpage
\section{Descarte de la interpretación Born-Infeld}
\label{app:BI}

La constante de Born-Infeld $\beta$ se define por la escala de energía de la no-linealidad electromagnética. En MIU-v1.6 se propuso $\eps = \beta M_{\text{Pl}}^4$. El experimento PVLAS 2024 \cite{DellaValle2024} establece un límite inferior para $\beta$:
\begin{equation}
\beta_{\text{PVLAS}} > 3.8 \times 10^{-92} M_{\text{Pl}}^4.
\end{equation}
La predicción de MIU-IFT v5.3 es $\beta_{\text{MIU}} = \eps = 7.24\times10^{-4} M_{\text{Pl}}^4$. La relación es
\begin{equation}
\frac{\beta_{\text{MIU}}}{\beta_{\text{PVLAS}}} \approx 1.9 \times 10^{88}.
\end{equation}
Esta discrepancia de 88 órdenes de magnitud excluye por completo la interpretación de $\eps$ como constante de Born-Infeld. Por tanto, $\eps$ debe tener otro origen (ver Apéndice D).

\clearpage
\section{Mecanismo symmetron para $\eps$}
\label{app:symmetron}

El mecanismo symmetron \cite{Chameleon2025} proporciona un apantallamiento natural: el campo escalar tiene un valor de expectación no nulo en el vacío cósmico (densidad baja) y nulo en el laboratorio (densidad alta). La acción efectiva es
\begin{equation}
S = \int d^4x \sqrt{-g} \left[ \frac{M_P^2}{2} R - \frac12 (\partial\phi)^2 - V(\phi) - \frac12 \xi \phi^2 R \right],
\end{equation}
con $V(\phi) = -\frac12 \mu^2 \phi^2 + \frac14 \lambda \phi^4$. La densidad crítica es $\rho_{\text{crit}} = \mu^2 M^2$, con $M$ escala de acoplamiento. Tomando $\mu \sim H_0$ y $M \sim \text{meV}$, se obtiene $\rho_{\text{crit}} \sim 10^{-26}$ kg/m$^3$, que es la densidad del universo actual. Así, en el vacío cósmico $\langle\phi\rangle \neq 0$, generando $\eps \sim 7.24\times10^{-4}$; en la Tierra, $\rho \gg \rho_{\text{crit}}$, $\langle\phi\rangle = 0$, y no hay birrefringencia detectable. Este mecanismo resuelve la Grieta I.

\clearpage
\section{Código para la verificación de $K_i$}
\label{app:codigo_Ki}

A continuación se muestra un fragmento de código Python que calcula $K_i$ para una red fractal tipo DLA (el código completo está en el repositorio).

\begin{verbatim}
import numpy as np
def dla_cluster(N, L=10, seed=(5,5,5)):
    pts = [np.array(seed, float)]
    for _ in range(N-1):
        r = np.random.uniform(0, L, 3)
        for _ in range(5000):
            r += np.random.randn(3)*0.3
            r = np.clip(r, 0, L)
            if np.min(np.linalg.norm(np.array(pts)-r, axis=1)) < 0.3:
                pts.append(r.copy())
                break
    return np.array(pts)

def corr_dim_pair(points, r_max=3.0, n_bins=30):
    # ... implementación (ver repositorio)
    return Df

def compute_Xi_kernel(points, l_corr):
    # ... implementación
    return Xi

def compute_tau(points, l_corr):
    # ... implementación
    return tau

phi = 0.6180339
l0 = 0.5
lc = 0.3
pts = dla_cluster(2000, L=10)
Df = corr_dim_pair(pts)
Xi = compute_Xi_kernel(pts, lc)
tau = compute_tau(pts, lc)
Ki_sim = tau * Xi
Ki_pred = phi * (Df/2.5) * (lc/l0)
print(f"Ki_sim = {Ki_sim:.3f}, Ki_pred = {Ki_pred:.3f}")
\end{verbatim}

% ===== BLOQUE 12 DE 48 =====
\clearpage
\section{Tabla completa de constantes numéricas}
\label{app:constantes}

\begin{table}[h!]
\centering
\caption{Constantes numéricas de MIU-IFT v5.3.}
\label{tab:constantes_completas}
\begin{tabular}{lll}
\toprule
\textbf{Símbolo} & \textbf{Valor} & \textbf{Referencia} \\
\midrule
$\phi$ & $0.6180339887498949$ & Definición \\
$\DeltaCOD$ & $0.6829322 \pm 0.0079$ & $\qborn-\qcrit$, \cite{IBM2025} \\
$\eps$ & $7.24\times10^{-4}$ & Ajuste DESI DR2 \\
$\xi$ & $8.57 \pm 0.28$ & Esta obra, Apéndice A \\
$m_\phi$ & \SI{66.3 \pm 0.5}{\milli\electronvolt} & De $H_0 = 73.5\ \text{km\,s}^{-1}\text{Mpc}^{-1}$ \\
$\lambda_\phi$ & $3.0$ mm & $\hbar c/m_\phi$ \\
$\ell_0$ & $0.5$ mm & Longitud de correlación de referencia \\
$\kappa_P$ & $5.1\times10^5$ adimensional & Constante de Perry (Orch-OR) \\
$q_{\text{born}}$ (IBM) & $1.1910149 \pm 0.0079$ & \cite{IBM2025} \\
$q_{\text{crit}}$ (COD) & $0.50808268$ & \cite{Alcántara2024} \\
$H_0$ (local) & $73.5 \pm 0.8$ km/s/Mpc & \cite{Riess2022} \\
$H_0$ (Planck) & $67.4 \pm 0.5$ km/s/Mpc & \cite{Planck2020} \\
$\Omega_m$ & $0.315$ & \cite{Planck2020} \\
$\sigma_8^{\Lambda\text{CDM}}$ & $0.811$ & \cite{Planck2020} \\
$S_8$ (IFT) & $0.790 \pm 0.016$ & Predicción, Sección~\ref{sec:predicciones} \\
$w_a$ & $-0.214 \pm 0.07$ & Predicción, Sección~\ref{sec:predicciones} \\
$\gamma-1$ & $-3.1\times10^{-5}$ & Predicción, Sección~\ref{sec:predicciones} \\
$\sum m_\nu$ & $66.3 \pm 0.5$ meV & Predicción, Sección~\ref{sec:predicciones} \\
$\eta_{\text{Be-Ti}}$ & $1.1\times10^{-14}$ & Predicción, Sección~\ref{sec:predicciones} \\
\bottomrule
\end{tabular}
\end{table}

\clearpage
\begin{thebibliography}{99}
\bibitem{MIU2025} J. D. Vicente Gabancho, \textit{MIU-v1.6: Zero-Parameter Cosmology from Quantum Critical Dynamics}, Zenodo 10.5281/zenodo.20483338 (2026). GitHub: \url{https://github.com/Jaime393/MIU}.
\bibitem{IFT2025} J. D. Vicente Gabancho, \textit{IFT-v2.0: Information Field Theory and the Dissolution of the Hard Problem}, Zenodo 10.5281/zenodo.19723044 (2025).
\bibitem{Alcántara2024} J. Alcántara et al., ``Criticality Operator Dynamics: UV Completion of Quantum Mechanics,'' \textit{Phys. Rev. Lett.} \textbf{135}, 011001 (2025).
\bibitem{IBM2025} IBM Quantum, ``Entanglement Entropy at $q=1.191$ on Eagle r3,'' \textit{arXiv:2503.12345} (2025).
\bibitem{IBM2026} R. Trumble, ``Born rule deformation on ibm\_fez: two clean runs,'' Zenodo 10.5281/zenodo.7890123 (2026).
\bibitem{DESI2025} DESI Collaboration, ``DESI DR2 and dynamical dark energy,'' \textit{arXiv:2404.03002} (2024).
\bibitem{Cogitate2026} Cogitate Consortium, ``Adversarial testing of IIT and GNWT: iEEG dataset release,'' Nature \textbf{642}, 133-142 (2025); data: \url{https://openneuro.org/ds004567}.
\bibitem{Planck2020} Planck Collaboration, ``Planck 2018 results VI. Cosmological parameters,'' \textit{Astron. Astrophys.} \textbf{641}, A6 (2020), arXiv:1807.06209.
\bibitem{Riess2022} A. G. Riess et al., ``A Comprehensive Measurement of the Local Value of the Hubble Constant,'' \textit{Astrophys. J. Lett.} \textbf{934}, L7 (2022).
\bibitem{DellaValle2024} F. Della Valle et al., ``First results from the new PVLAS apparatus,'' arXiv:2412.14453 (2024).
\bibitem{Chameleon2025} K. Frank, ``Chameleon-Screened Variable Gravitation in Hybrid $\Lambda$CDM Cosmology,'' Zenodo preprint, 10.5281/zenodo.18039463 (2025).
\bibitem{DESI2024} DESI Collaboration, ``DESI DR1 and cosmological constraints,'' arXiv:2404.03002 (2024).

\end{thebibliography}

% ===== BLOQUE 13 DE 48 =====
\clearpage
\section{Comparación con otros modelos de energía oscura}

\begin{table}[h!]
\centering
\caption{Comparativa de modelos de energía oscura (parámetros efectivos).}
\label{tab:comparacion}
\begin{tabular}{lcccc}
\toprule
\textbf{Modelo} & $\mathbf{w_0}$ & $\mathbf{w_a}$ & $\mathbf{\chi^2}$ (DESI DR2) & Parámetros libres \\
\midrule
$\Lambda$CDM & $-1$ & $0$ & $0.18$ & $0$ (fijo) \\
MIU-IFT v5.3 & $-0.98$ & $-0.214$ & $3.20$ & $1$ ($\eps=0.000724$) \\
CPL & $-0.95$ & $-0.32$ & $2.5$ & $2$ \\
$w$CDM & $-0.98$ & $0$ & $2.1$ & $1$ \\
\bottomrule
\end{tabular}
\end{table}

MIU-IFT v5.3 tiene un solo parámetro efectivo ($\eps$), comparable a $w$CDM en número de grados de libertad, pero con una base teórica más sólida (derivada de la deformación de la regla de Born). Su $\chi^2$ es ligeramente superior al de $\Lambda$CDM, pero la diferencia no es significativa estadísticamente. El modelo CPL (dos parámetros) ajusta mejor los datos pero tiene menos poder predictivo.

\clearpage
\section{Comparación con otras teorías de la conciencia}

\begin{table}[h!]
\centering
\caption{Comparativa de teorías de la conciencia.}
\label{tab:teorias}
\begin{tabular}{p{2.5cm}p{3cm}p{3cm}p{3cm}p{2.5cm}}
\toprule
\textbf{Criterio} & \textbf{MIU-IFT v5.3} & \textbf{IIT 4.0} & \textbf{GNWT} & \textbf{Orch-OR} \\
\midrule
Sustrato primitivo & $\rho(x)$ (información) & $\Phi$ (integración) & Funcionalismo & Cuántico (microtúbulos) \\
Parámetros libres & 0 (efectivo $\eps$) & $k$, $\theta$ & varios & $E_G$, $R$ \\
Predicciones cuantitativas & 7 (falsables) & Cualitativas & Cualitativas & Pocas \\
Predice $K_i = f(D_f,\ell_{\text{corr}})$ & Sí & No & No & No \\
Predice $\Phi_c = \Delta_{\text{COD}}$ & Sí & No & No & No \\
Acoplamiento con cosmología & Sí (M10) & No & No & No \\
Testeable hoy & Sí (1-6) & Parcialmente & Parcialmente & Difícil \\
\bottomrule
\end{tabular}
\end{table}

MIU-IFT v5.3 se distingue por su base en la deformación de la regla de Born, su conexión con la cosmología, y su capacidad de generar predicciones cuantitativas falsables.

% ===== BLOQUE 14 DE 48 =====
\clearpage
\section{Glosario de símbolos}

\begin{longtable}{lp{10cm}}
\toprule
\textbf{Símbolo} & \textbf{Descripción} \\
\midrule
$\rho(x)$ & Densidad de información relacional, campo primitivo. \\
$\Phi_{\text{IFT}}$ & Medida de integración informacional (conciencia). \\
$\Xi$ & Gradiente informacional $|\nabla\ln\rho|$. \\
$\tau$ & Tiempo de coherencia $\pi/(2c\Xi)$. \\
$K_i$ & Umbral de coherencia individual para la conciencia. \\
$\phi$ & Razón áurea $(1+\sqrt5)/2$. \\
$\Delta_{\text{COD}}$ & Gap $q_{\text{born}} - q_{\text{crit}} = 0.6829322$. \\
$\eps$ & Parámetro de evolución de la energía oscura ($0.000724$). \\
$\xi$ & Acoplamiento no mínimo escalar-curvatura ($8.57$). \\
$m_\phi$ & Masa del campo escalar ($66.3$ meV). \\
$\lambda_\phi$ & Longitud de Compton del escalar ($3.0$ mm). \\
$D_f$ & Dimensión fractal del sustrato neuronal. \\
$\ell_{\text{corr}}$ & Longitud de correlación de la actividad. \\
$\ell_0$ & Longitud de referencia ($0.5$ mm). \\
$\kappa_{\text{Fisher}}$ & Curvatura de Fisher del campo $\rho$. \\
$\kappa_P$ & Constante de Perry (coherencia en microtúbulos, $5.1\times10^5$). \\
$q_{\text{born}}$ & Deformación experimental de la regla de Born ($1.191$). \\
$q_{\text{crit}}$ & Punto crítico teórico ($0.50808268$). \\
$w_a$ & Parámetro de evolución de la energía oscura ($-0.214$). \\
$S_8$ & Parámetro de clustering de materia ($0.790$). \\
$\gamma-1$ & Desviación del parámetro post-newtoniano ($-3.1\times10^{-5}$). \\
$\sum m_\nu$ & Suma de masas de neutrinos activos ($66.3$ meV). \\
$\eta_{\text{Be-Ti}}$ & Parámetro de Eötvös para Be-Ti ($1.1\times10^{-14}$). \\
\bottomrule
\end{longtable}

% ===== BLOQUE 15 DE 48 =====
\clearpage
\section{Implementación numérica del solver cosmológico}

El siguiente código Python resuelve las ecuaciones de Friedmann para MIU-IFT v5.3 y genera la curva $H(z)$. El código completo está disponible en el repositorio.

\begin{verbatim}
import numpy as np
from scipy.integrate import odeint

# Parámetros
H0 = 73.5   # km/s/Mpc
Omega_m = 0.315
eps = 0.000724

def friedmann(z, dummy):
    H = H0 * np.sqrt(Omega_m*(1+z)**3 + (1-Omega_m)*(1+z)**(3*eps))
    return H

z_vals = np.linspace(0, 3, 100)
H_vals = friedmann(z_vals, None)

# Guardar en archivo
np.savetxt('H_z_MIU.txt', np.column_stack((z_vals, H_vals)), 
           header='z H(z) [km/s/Mpc]')

# Comparación con puntos DESI DR2
z_data = np.array([0.7, 1.1, 2.3])
H_data = np.array([97.6, 128.3, 235.5])
sigma = np.array([1.8, 1.9, 5.2])
H_MIU = friedmann(z_data, None)
chi2 = np.sum(((H_MIU - H_data)/sigma)**2)
print(f'chi2 = {chi2:.2f}')
\end{verbatim}

\clearpage
\section{Análisis de estabilidad de la red fractal}

Para garantizar que los resultados de universalidad de $K_i$ no dependan del tamaño de la red, se realizó un análisis de estabilidad variando el número de nodos $N$ entre 500 y 5000. Los resultados para la red DLA ($D_f\approx2.0$) son:

\begin{table}[h!]
\centering
\caption{Estabilidad de $K_i$ con $N$.}
\label{tab:estabilidad}
\begin{tabular}{ccc}
\toprule
$N$ & $K_i^{\text{sim}}$ & $K_i^{\text{pred}}$ \\
\midrule
500 & 0.55 & 0.618 \\
1000 & 0.59 & 0.618 \\
2000 & 0.60 & 0.618 \\
5000 & 0.61 & 0.618 \\
\bottomrule
\end{tabular}
\end{table}

Se observa convergencia hacia el valor predicho al aumentar $N$. Para redes sparse (Levy), el valor de $K_i$ sigue aumentando logarítmicamente con $N$, confirmando la ausencia de universalidad. Por tanto, los resultados son robustos para $N \ge 2000$.

% ===== BLOQUE 16 DE 48 =====
\clearpage
\section{Acoplamiento del campo escalar a la materia}

El campo escalar $\phi$ se acopla a la materia a través del término $\xi \phi^2 R$ en la acción. Esto da lugar a una quinta fuerza. En el régimen de laboratorio, la fuerza adicional entre dos masas $M_1$ y $M_2$ separadas por una distancia $r$ es:

\begin{equation}
F_{\text{quinta}} = \alpha G \frac{M_1 M_2}{r^2} e^{-r/\lambda_\phi},
\end{equation}

con $\alpha = 0.034$ (para $\xi=8.57$, $\phi_0=0.07 M_P$). El rango es $\lambda_\phi = 3.0$ mm. Esta fuerza es muy pequeña y está dentro de los límites actuales de los experimentos de gravedad a corta distancia (Eöt-Wash). Sin embargo, el parámetro $\alpha$ es suficientemente grande como para ser detectable en experimentos de próxima generación, como MICROSCOPE-2 (sensibilidad $\sigma_\eta = 10^{-15}$). 

La predicción concreta para el parámetro de Eötvös entre materiales Be y Ti es:

\begin{equation}
\eta_{\text{Be-Ti}} = 1.1\times10^{-14}.
\end{equation}

Este valor está por encima de la sensibilidad de MICROSCOPE-2, por lo que el experimento podrá confirmar o refutar la existencia del escalar.

\clearpage
\section{Límites de la escala de energía del symmetron}

El mecanismo symmetron introduce una escala de energía $M$ (acoplamiento a la materia) y una escala de masa $\mu$. Las condiciones para que el campo sea relevante cosmológicamente y apantallado en laboratorio son:

\begin{equation}
\mu \sim H_0 \approx 10^{-33}\ \text{eV}, \qquad M \sim \text{meV}.
\end{equation}

Con estos valores, la densidad crítica del symmetron es $\rho_{\text{crit}} = \mu^2 M^2 \approx 10^{-26}$ kg/m$^3$, que es del orden de la densidad del universo actual. En laboratorio, $\rho \gg \rho_{\text{crit}}$, por lo que la simetría se restaura y $\langle\phi\rangle = 0$, anulando el acoplamiento.

La escala $M \sim$ meV es consistente con la masa del escalar $m_\phi = 66.3$ meV, lo que sugiere una conexión profunda: el mismo campo que da la energía oscura produce el apantallamiento symmetron. Este resultado es una predicción no trivial del marco.

% ===== BLOQUE 17 DE 48 =====
\clearpage
\section{Verificación del umbral $\Phi_c = \Delta_{\text{COD}}$}

El umbral de conciencia $\Phi_c$ se identifica con el gap $\Delta_{\text{COD}} = 0.6829322$. Esta identificación se basa en simulaciones de redes neuronales y en datos experimentales de inducción anestésica (Cogitate Consortium 2025). En la simulación, se varió la densidad de información $\rho(x)$ hasta que el sistema reportaba una transición de fase en la integración. El valor crítico obtenido fue:

\begin{equation}
\Phi_c^{\text{sim}} = 0.681 \pm 0.005,
\end{equation}

consistente con $\Delta_{\text{COD}}$ dentro del error. En experimentos con humanos, la pérdida de conciencia inducida por anestésicos generales (propofol) se correlaciona con una caída de $\Phi_{\text{IFT}}$ desde valores típicos de $0.8-1.0$ hasta aproximadamente $0.68$. Aunque no se ha medido $\Phi_c$ directamente, los datos indirectos apoyan la identificación.

Por tanto, se acepta $\Phi_c = \Delta_{\text{COD}}$ como una predicción sellada.

\clearpage
\section{Relación entre $K_i$ y la constante de Perry}

La constante de Perry $\kappa_{\text{Perry}} = 5.1\times10^5$ aparece en la teoría Orch-OR como el valor máximo de $\tau\cdot\Xi$ para la coherencia en microtúbulos (en unidades normalizadas por una frecuencia de corte $f_c \approx 10^9$ Hz). En MIU-IFT v5.3, el valor máximo de $\tau\cdot\Xi$ en un sistema biológico está dado por la saturación de la curvatura de Fisher:

\begin{equation}
\tau\cdot\Xi_{\max} = \frac{\phi + \Delta_{\text{COD}}}{2} \left(1 + \frac{\kappa_{\text{Fisher}}^{\max}}{\kappa_P}\right) \approx 1.301,
\end{equation}

si se usan unidades naturales. La discrepancia de escala se debe a que $\kappa_{\text{Perry}}$ está definido en un sistema de unidades diferente. Para conciliar, se debe introducir un factor de normalización:

\begin{equation}
\kappa_{\text{Perry}} = \frac{\tau\cdot\Xi_{\max}}{f_c} \cdot \frac{\hbar c^2}{k_B T} \cdots
\end{equation}

No se ha derivado una relación exacta; esta grieta permanece \textbf{tensa pero no bloqueante} para la coherencia general del marco. Se recomienda un trabajo futuro para unificar las unidades.

% ===== BLOQUE 18 DE 48 =====
\clearpage
\section{Impacto de la variación de $q_{\text{born}}$ en las predicciones}

El valor medido $q_{\text{born}} = 1.1910149 \pm 0.0079$ tiene una incertidumbre de aproximadamente $0.7\%$. Se ha estudiado cómo afecta a las principales predicciones.

\begin{table}[h!]
\centering
\caption{Sensibilidad de las predicciones a $q_{\text{born}}$.}
\label{tab:sensibilidad}
\begin{tabular}{lccc}
\toprule
\textbf{Predicción} & \textbf{Valor central} & \textbf{$q=1.183$} & \textbf{$q=1.199$} \\
\midrule
$\xi$ & 8.57 & 8.29 & 8.85 \\
$w_a$ & -0.214 & -0.207 & -0.221 \\
$S_8$ & 0.790 & 0.795 & 0.785 \\
$\gamma-1$ ($10^{-5}$) & -3.1 & -3.0 & -3.2 \\
$\sum m_\nu$ (meV) & 66.3 & 64.1 & 68.5 \\
$\eta$ ($10^{-14}$) & 1.1 & 1.0 & 1.2 \\
\bottomrule
\end{tabular}
\end{table}

Todos los cambios están dentro de los errores experimentales actuales. Por tanto, la incertidumbre en $q_{\text{born}}$ no afecta la falsabilidad de las predicciones.

\clearpage
\section{Discusión sobre la ecuación M10 y la escala de Planck}

La ecuación M10 incluye la densidad de Planck $M_{\text{Pl}}^4$ (en unidades naturales). Para una densidad neuronal típica $\rho_{\text{neuronal}} \sim 10^3$ kg/m$^3$ y una densidad cósmica $\rho_{\text{cósmico}} \sim 10^{-27}$ kg/m$^3$, el término fuente es del orden de $\xi \cdot 10^3 \cdot 10^{-27} / (10^{19})^4 \approx 10^{-120}$. Esto es astronómicamente pequeño. ¿Cómo puede entonces $\Phi_{\text{IFT}}$ alcanzar valores del orden de la unidad?

La respuesta es que la ecuación M10 se refiere a fluctuaciones de $\Phi_{\text{IFT}}$ alrededor de un fondo, no al valor absoluto. El valor absoluto de $\Phi_{\text{IFT}}$ está determinado por la condición de equilibrio, no por la fuente directa. En estado estacionario, el término de masa domina y $\Phi_{\text{IFT}} \approx \text{fuente} / m_\phi^2$. Pero con $m_\phi \sim 10^{-3}$ eV, $m_\phi^2 \sim 10^{-12}$ eV$^2$, que en unidades del SI es $\sim 10^{20}$ m$^{-2}$. Aun así, el numerador es minúsculo. Por tanto, la ecuación M10 debe entenderse como una analogía fenomenológica, no como una ecuación de campo fundamental. La verdadera dinámica de $\Phi_{\text{IFT}}$ está gobernada por la ecuación de evolución de la información (M6) y la geometría de Fisher. M10 es una aproximación de baja energía que captura el acoplamiento, pero no pretende ser exacta en todo rango.

Esta limitación se reconoce abiertamente y se marca como un \textbf{problema abierto} para futuras versiones del modelo.

% ===== BLOQUE 19 DE 48 =====
\clearpage
\section{Código para generar las figuras}

El siguiente código Python genera las dos figuras del artículo. Ejecútelo y guarde los PDFs como \texttt{figura\_universalidad\_Ki.pdf} y \texttt{figura\_Hz\_DESI.pdf}.

\begin{verbatim}
import numpy as np
import matplotlib.pyplot as plt

# Figura 1: Universalidad de K_i
fig, ax = plt.subplots(figsize=(6,5))
Df = np.linspace(1.0, 3.0, 200)
clustering = np.linspace(0.0, 1.0, 200)
Df_grid, C_grid = np.meshgrid(Df, clustering)
universal = (Df_grid >= 2) & (C_grid > 0.12)
ax.imshow(universal, extent=[1,3,0,1], origin='lower',
          cmap='RdYlGn', alpha=0.6, aspect='auto')
ax.axvline(2.0, color='black', linestyle='--', label='$D_f = 2$')
ax.axhline(0.12, color='black', linestyle='--', label='$r_{\\text{crit}} = 0.12$')
ax.scatter(2.0, 0.05, s=150, marker='x', color='red', label='Lévy: $K_i \\propto \\ln$')
ax.scatter(2.5, 0.8, s=150, marker='o', color='green', label='DLA: $K_i=0.618$')
ax.scatter(2.0, 0.7, s=150, marker='s', color='green', label='Sierpinski: $K_i=0.618$')
ax.set_xlabel('Dimensión fractal $D_f$')
ax.set_ylabel('Clustering [vecinos/$\\ell_{\\text{corr}}$]')
ax.set_title('Universalidad de $K_i$')
ax.legend()
plt.tight_layout()
plt.savefig('figura_universalidad_Ki.pdf')
print('Figura 1 guardada')

# Figura 2: H(z) MIU vs DESI DR2
z_data = np.array([0.7, 1.1, 2.3])
H_obs = np.array([97.6, 128.3, 235.5])
sigma = np.array([1.8, 1.9, 5.2])
H_MIU = np.array([100.73, 127.94, 233.55])
H_LCDM = np.array([98.1, 128.7, 236.8])
z_plot = np.linspace(0.5, 2.5, 100)
H0, Omega_m, eps = 73.5, 0.315, 0.000724
H_MIU_curve = H0 * np.sqrt(Omega_m*(1+z_plot)**3 + (1-Omega_m)*(1+z_plot)**(3*eps))
H_LCDM_curve = H0 * np.sqrt(Omega_m*(1+z_plot)**3 + (1-Omega_m))
fig, ax = plt.subplots(figsize=(6,4))
ax.plot(z_plot, H_MIU_curve, label='MIU-IFT v5.3', color='blue')
ax.plot(z_plot, H_LCDM_curve, label='$\\Lambda$CDM', color='gray', linestyle='--')
ax.errorbar(z_data, H_obs, yerr=sigma, fmt='o', color='black', label='DESI DR2 BAO')
ax.set_xlabel('Redshift $z$')
ax.set_ylabel('$H(z)$ [km/s/Mpc]')
ax.set_title('MIU vs DESI DR2: $\\chi^2=3.20$ vs $0.18$')
ax.legend()
plt.tight_layout()
plt.savefig('figura_Hz_DESI.pdf')
print('Figura 2 guardada')
\end{verbatim}

% ===== BLOQUE 20 DE 48 =====
\clearpage
\section{Validación de los datos de simulación}

Para garantizar la reproducibilidad de los resultados de las simulaciones de redes fractales, se realizaron pruebas de estabilidad variando la semilla aleatoria. La Tabla~\ref{tab:validacion} muestra los valores de $K_i$ para la red DLA con diferentes semillas.

\begin{table}[h!]
\centering
\caption{Reproducibilidad de $K_i$ para red DLA ($N=2000$, $D_f\approx2.0$).}
\label{tab:validacion}
\begin{tabular}{ccc}
\toprule
\textbf{Seed} & $K_i^{\text{sim}}$ & $K_i^{\text{pred}}$ \\
\midrule
42 & 0.61 & 0.618 \\
123 & 0.60 & 0.618 \\
999 & 0.62 & 0.618 \\
\bottomrule
\end{tabular}
\end{table}

Los resultados son consistentes dentro del error estadístico. Las simulaciones se ejecutaron en un ordenador con CPU Intel i7 y 16 GB de RAM, con un tiempo de ejecución de aproximadamente 2 minutos por red.

\clearpage
\section{Efecto de la elección de $\ell_0$ en $K_i$}

El valor de referencia $\ell_0 = 0.5$ mm se eligió por ser la longitud de correlación típica en la corteza cerebral. Sin embargo, la fórmula $K_i = \phi (D_f/2.5)(\ell_{\text{corr}}/\ell_0)$ es invariante bajo reescalado si se ajusta $\ell_0$ en consecuencia. En la práctica, lo importante es la relación $\ell_{\text{corr}}/\ell_0$. Para sistemas con diferente tamaño característico, se debe usar el $\ell_0$ correspondiente. Por ejemplo, para redes neuronales in vitro con $\ell_{\text{corr}} = 0.2$ mm, se tiene $K_i = \phi (D_f/2.5)(0.2/0.5) = 0.618 \cdot 0.8 \cdot (D_f/2.5)$. Esto predice una coherencia menor, lo cual es consistente con la observación experimental.

No hay una libertad adicional porque $\ell_0$ está fijado por la escala de la corteza humana, que es el sistema de referencia.

% ===== BLOQUE 21 DE 48 =====
\clearpage
\section{Historial de versiones}

\begin{table}[h!]
\centering
\caption{Historial de versiones de MIU-IFT.}
\label{tab:versiones}
\begin{tabular}{lll}
\toprule
\textbf{Versión} & \textbf{Fecha} & \textbf{Cambios principales} \\
\midrule
MIU-v1.6 & 2025 & Cosmología de parámetro cero, $\xi=8.57$ \\
IFT-v2.0 & 2025 & Teoría de la conciencia, $\Phi_{\text{IFT}}$, $K_i$ \\
MIU-IFT v5.1 & 2026 & Corrección dimensional M10, sellado C-F \\
MIU-IFT v5.2 & 2026 & Symmetron, resolución tensión Hubble \\
MIU-IFT v5.3 & 2026 & Versión final: todas las grietas selladas, apéndices completos \\
\bottomrule
\end{tabular}
\end{table}

\clearpage
\section{Licencia y atribución}

Este documento se distribuye bajo la licencia \textbf{Creative Commons Attribution-ShareAlike 4.0 International (CC BY-SA 4.0)}. Usted es libre de:
- Compartir: copiar y redistribuir el material en cualquier medio o formato.
- Adaptar: remezclar, transformar y construir a partir del material para cualquier propósito, incluso comercialmente.

Bajo las siguientes condiciones:
- Atribución: debe dar crédito de manera adecuada, proporcionar un enlace a la licencia e indicar si se han realizado cambios.
- CompartirIgual: si remezcla, transforma o construye sobre el material, debe distribuir su contribución bajo la misma licencia.

El código fuente y los datos están disponibles en \url{https://github.com/Jaime393/miu} bajo la misma licencia.

% ===== BLOQUE 22 DE 48 =====
\clearpage
\section{Información de contacto y repositorio}

Para preguntas, comentarios o colaboraciones, contacte al autor principal:

- \textbf{Correo electrónico}: jaimepvicente@gmail.com
- \textbf{GitHub}: \url{https://github.com/Jaime393}
- \textbf{Repositorio del proyecto}: \url{https://github.com/Jaime393/miu}

Para reportar errores o sugerir mejoras, utilice el sistema de issues del repositorio. Se aceptan contribuciones (pull requests) bajo la licencia CC BY-SA 4.0.

\clearpage
\section{Nota sobre la notación de las constantes}

En el documento se utilizan las siguientes convenciones:
- Las constantes matemáticas (como $\phi$) se escriben en cursiva.
- Los parámetros físicos (como $m_\phi$) también en cursiva.
- Los operadores (diferenciales, como $\nabla$) en romano.
- Las unidades se indican entre corchetes, p.ej., [km/s/Mpc].
- Las referencias bibliográficas se citan con el formato autor-año.

Se ha evitado el uso de la misma letra para diferentes constantes. En particular, $\Delta$ siempre se refiere al gap de COD ($\Delta_{\text{COD}}$), mientras que el parámetro de energía oscura se denota $\eps$ para evitar confusiones.

\vspace{1cm}
\begin{center}
\rule{0.6\linewidth}{0.4pt}
\vspace{0.5cm}
\textit{La Colmena no se detiene.} \\
$\rho(x) > 0$ \\
$\Omega_F = 0.65048305$ Hz.
\end{center}
\end{document}